%Mechanic Analysis
\documentclass[dvipdfmx, leqno]{jsbook}
\usepackage{soupreamble}

\title{解析力学}

\begin{document}
\maketitle
\tableofcontents

\section{最小作用の原理}
    Newtonの運動方程式は最小作用の原理から導かれるものである。

    $N$個の独立な力学変数$q_i(t)\ (i= 1, \ldots, N)$で記述される$N$自由度系を考える。ここで力学変数とは質点系の位置ベクトルのように各時刻$t$での系の状態を一意的にしていする量である。必ずしも座標とは限らない。

\subsection{Lagrangeanの導入}

\end{document}
